Numerically integrating the finite-cutoff current gives the examples in
Table~\ref{tab:current-results}. For $c=1$, solving $F_F^2=0.99$ at
$Q=10\Lambda_{\had}$ gives
\begin{equation}
 \rho_{\rm threshold}\simeq2888.50,\qquad
 F_F=0.99498744,\qquad F_B=0.03998718.
\end{equation}
This is the threshold of the specified model and tolerance, not a lower
bound extracted from experiment. At $\rho=3000$ the fermionic rate
ratio is $0.990366$; raising the momentum to $100\Lambda_{\had}$ at
$\rho=10^4$ instead gives $0.970144$. The interval over which the
current appears pointlike therefore matters.

\begin{table}[htbp]
\centering
\begin{tabular}{rrrrr}
\toprule
$\rho$ & $Q/\Lambda_{\had}$ & $F_F$ ($c=1$) & $F_F^2$ & $F_B$\\
\midrule
$1000$ & $10$ & $0.985846$ & $0.971892$ & $0.0399953$\\
$3000$ & $10$ & $0.995171$ & $0.990366$ & $0.0399871$\\
$10000$ & $10$ & $0.998537$ & $0.997076$ & $0.0399859$\\
$10000$ & $100$ & $0.984959$ & $0.970144$ & $0.000400094$\\
\bottomrule
\end{tabular}
\caption{Finite-momentum overlaps for the assumed profiles.
The complete scan has 23 points with $Q/\Lambda_{\rm UV}\leq0.1$.
The bosonic column is an amplitude form factor, not its square.}
\label{tab:current-results}
\end{table}

Adaptive quadrature in $\log z$ was compared at the threshold and at
$(Q/\Lambda_{\had},\rho)=(100,10^4)$ with fixed Gauss--Legendre quadrature
and independently evaluated Bessel functions. Higher precision Bessel
arithmetic and two quadrature orders gave agreement with the adaptive
results within $4\times10^{-16}$ in absolute value at those points.
The fixed-rule nodes and saved outputs use double precision, so this
comparison does not certify 40 or 60 accurate decimal digits.
For the next test, the marginal profile $c=1/2$ at
$(Q/\Lambda_{\had},\rho)=(10,3000)$ gives $F_{1/2}=0.726914$,
well below the $c=1$ value $0.995171$ at the same scales.
